\documentclass[12pt]{article}

\usepackage{sbc-template}
\usepackage{graphicx,url}
\usepackage[utf8]{inputenc}
\usepackage[brazil]{babel}
%\usepackage[latin1]{inputenc}  

     
\sloppy

\title{Relatório de Projetos de Engenharia 1\\ FeedPet Comedouro Automático}

\author{Emily D. Guimarães\inst{1}, Flávio S. Martins Filho\inst{1}, Lauro V. Garcês Lima\inst{1},\\
Marcos Lobato Chagas\inst{1}, Sophia A. de Sousa da Silva\inst{1}}


\address{Instituto de Tecnologia -- Universidade Federal do Pará \\
  Caixa Postal 479 – 66075-110 – Belém – Pará – Brasil \\
  \email{
    \begin{tabular}{c}
        \{emily.guimaraes, flavio.filho, lauro.lima, marcos.chagas, \\
        sophia.silva\}@itec.ufpa.br
    \end{tabular}
    }
}



\begin{document} 

\maketitle

\begin{abstract}

  This report details the development of an automated and intelligent pet feeding and watering system integrated with an Android smartphone application. The system aims to solve common modern challenges in pet care, ensuring precise nutritional portioning and fresh water availability. Utilizing an ESP32 microcontroller, load cells, ultrasonic sensors, and solenoid valves, the hardware interacts seamlessly with a mobile application via Wi-Fi. The article details the architecture, methodology, implementation, and verification processes of the prototype.
  
\end{abstract}
     
\begin{resumo}

  Este relatório detalha o desenvolvimento de um sistema automatizado e inteligente de alimentação e hidratação para animais de estimação integrado a um aplicativo para smartphone Android. O sistema visa resolver desafios modernos no cuidado de pets, garantindo porcionamento nutricional preciso e disponibilidade de água fresca. Utilizando um microcontrolador ESP32, células de carga, sensores ultrassônicos e válvulas solenoides, o hardware interage diretamente com um aplicativo móvel via Wi-Fi. O artigo detalha a arquitetura, metodologia, implementação e os processos de verificação do protótipo.

\end{resumo}

\section{Introdução}

O projeto propõe formular um alimentador automático integrado com um aplicativo para smartphones Android que atenda às necessidades dos donos de cães e gatos.

Ultimamente, é possível observar que diversos tutores estão procurando certas facilidades no seu dia a dia e também produtos que garantem qualidade de vida para os animais domésticos. Em contrapartida, de acordo com o \cite{folha:23}, cerca de 25\% a 40\% dos animais no Brasil são obesos, isso significa que o cuidado alimentar não está sendo devidamente controlado.

Para garantir que o projeto seja eficiente, o diferencial é a integração do aplicativo no dispositivo. Dessa forma, o tutor do animal poderá ter praticidade no dia a dia, visto que a responsabilidade de controlar o cronograma alimentar do animal é passada para o alimentador eletrônico. Além disso, por contra desse monitoramento, ele é ideal para rotinas imprevisíveis, uma vez que o dispositivo será controlado por uma dosagem correta. Assim, os donos poderão monitorar a rotina alimentar de seus animais. Por fim, a união desses fatores garante a disciplina e a qualidade de vida para esses cães e gatos.

O objetivo do projeto é criar um comedouro e bebedouro inteligente com o auxílio de um aplicativo, a fim de facilitar o cuidado do tutor com o animal e também garantir qualidade de vida para cães e gatos. O dispositivo é focado em animais de pequeno a médio porte, contendo um pote para ração e um para a água. Já o aplicativo é responsável por selecionar os intervalos e horários que o alimento e a água serão despejados no pote do dispositivo.

\section{Trabalhos Relacionados}

Para a formulação de um sistema embarcado que atenda a todos os requisitos até então expostos, é necessário antes analisar as opções já comercializadas.

\subsection{Soluções existentes}

As soluções já adotadas no mercado, como o \cite{onizoomi:26} e o \cite{surecare:26}, possuem limitações simples de serem resolvidas, a saber: o alto custo monetário, a venda separada do comedouro e bebedouro (ou a ausência completa do bebedouro) e a baixa capacidade de personalização da ração e água disponibilizada. Isso se deve ao fato desses produtos priorizarem a praticidade acima da individualidade de cada animal doméstico, o que ocasiona na padronização de um sistema universal que se limita à porções de ração pré-definidas com o aplicativo - o qual não está incluso em alguns dos produtos disponíveis no mercado atualmente.

\section{Arquitetura do Sistema / Metodologia}

\subsection{Proposta}

Com isso, a ideia inicial do projeto foi solucionar essas limitações dos produtos à venda no mercado atual. Para que o objetivo seja atingido, foi necessário adotar duas medidas imprescindíveis: a criação de um dispositivo único que gerenciasse as porções de ração e a quantidade de água despejada e a formulação  de um app com baixo custo de manutenção. 

\subsection{Dispositivo}



\subsection{Aplicativo}



\subsection{Fluxogramas}



\section{Avaliação e Resultados}



\section{Custos do projeto}

\begin{table}[ht]
    \centering
    \caption{Custos do projeto}
    \label{tab:consumo_esp32}
    \vskip 0.1in % Garante o espaçamento correto exigido pela SBC abaixo da legenda
    \begin{tabular}{|l|c|r|}
    \hline
    \textbf{Componente}        		           		& \textbf{Quantidade}    	& \textbf{Preço}    		\\ \hline
    Microcontrolador ESP32    			       	& 1                      			& R\$40,90          		\\ \hline
    Módulo L293N                    	       			& 1                      			& R\$19,50          		\\ \hline
    Sensor de nível e detecção de água         		& 1            		    		& R\$08,60          		\\ \hline
    Célula de carga + 1 Módulo Hx711	       		& 1						& R\$11,30			\\ \hline
    Motor 9V								& 1						& R\$18,90			\\ \hline
    Bomba 9V								& 1						& R\$24,00			\\ \hline
    Buck Boost							       	& 1						& R\$12,00			\\ \hline
    Bateria 9V							       	& 1						& R\$				\\ \hline
    LED									& 						& R\$02,00			\\ \hline
    Botão								       	& 2						& R\$03,00			\\ \hline
    Buzzer							           	& 1						& R\$05,00			\\ \hline
    Protoboard							       	& 1						& R\$10,00			\\ \hline
    Jumpers								& 						& R\$10,00		    	\\ \hline
    Resistores							       	& 						& R\$05,00			\\ \hline
    \multicolumn{2}{|l|}{\textbf{Total}}       		& \textbf{R\$170,00}				            		\\ \hline
    \end{tabular}
\end{table}

\section{Conclusão}

A construção do dispositivo e do aplicativo foram criados, respectivamente, com o auxílio de ferramentas onlines, como \cite{mermaid:26}, \cite{tinkercad:23}, \cite{gettyimages:19}, \cite{blueoprint:26} e \cite{pandasuite:26}. Já as soluções existentes disponíveis nos sites \cite{onizoomi:26}, \cite{surecare:26}, \cite{freepik:26} serviram de base para a consolidação da proposta do projeto. Ademais, os  seguintes dados \cite{folha:23} contribuiram para fundamentar a visão do projeto. 

\bibliographystyle{sbc}
\bibliography{sbc-template}

\end{document}
